% 脉冲星行星（综述）
% license CCBYSA3
% type Wiki

本文根据 CC-BY-SA 协议转载翻译自维基百科\href{https://en.wikipedia.org/wiki/Pulsar_planet}{相关文章}。

\begin{figure}[ht]
\centering
\includegraphics[width=6cm]{./figures/427823341c42e767.png}
\caption{带有行星的脉冲星艺术概念图} \label{fig_Pupl_1}
\end{figure}
脉冲星行星是指围绕脉冲星运行的行星。1992 年，人们首次在一颗毫秒脉冲星周围发现了这类行星，它们也是人类确认发现的第一批太阳系外行星。脉冲星本身是极其精确的“宇宙时钟”，即使质量很小的行星，也能在脉冲星的观测特性中引起可探测的变化；目前已知质量最小的系外行星，正是一颗脉冲星行星。

这类行星极为罕见，NASA 系外行星档案中仅列出大约半打（六颗左右）。只有一些非常特殊的过程，才能在脉冲星周围形成行星级别的伴星，其中许多被认为是性质十分奇特的天体，例如由伴星部分被破坏后形成的“钻石行星”。此外，脉冲星强烈的辐射以及由电子–正电子对组成的高速粒子风，往往会剥离这些行星的大气层，因此它们几乎不可能成为生命适宜的栖居地。
\subsection{形成}
行星的形成需要原行星盘的存在，而多数理论还要求盘内存在一个“死区”，即没有湍流的区域。在该区域内，微行星体可以形成并不断积累，而不会因向恒星内落而被吞噬$^{[1]}$。与年轻恒星相比，脉冲星具有高得多的光度，因此其辐射会强烈电离盘物质，从而抑制死区的形成$^{[2]}$；电离会触发磁旋转不稳定性，引发湍流并破坏死区$^{[3]}$。因此，若要孕育行星，围绕脉冲星的盘必须具有相当大的质量$^{[4]}$。
目前提出了若干可能形成行星系统的过程$^{[a]}$：
\begin{itemize}
\item “第一代”行星：指在恒星发生超新星爆发并演化为脉冲星之前，就已围绕其运行的行星$^{[6]}$。大质量恒星往往缺乏行星，这可能既源于在极其明亮的恒星周围难以探测行星，也由于强辐射会破坏原行星盘。若行星轨道半径小于约 4 个天文单位，当恒星演化为红巨星或红超巨星时，行星很可能被吞没并毁灭。在超新星爆发过程中，系统会损失大约一半的质量；除非脉冲星在爆发时被“踢出”的方向恰与行星当时的运动方向一致，否则行星往往会与系统解体而脱离。已知的脉冲星行星系统中，没有一个很可能是通过这一过程形成的$^{[7]}$。
\item “第二代”行星：由超新星爆发后回落到脉冲星上的物质形成$^{[6]}$。理论上，这些回落物质的总质量可与典型原行星盘相当$^{[7]}$，但其耗散速度可能过快，不利于行星形成。目前尚无围绕年轻脉冲星的已知行星实例$^{[8][9]}$。
\item “第三代”行星$^{[6]}$：脉冲星与伴星相互作用并将其摧毁，从而形成一个低质量盘。脉冲星可发出高能辐射加热伴星，使其充满并溢出罗希瓣，最终被完全破坏。另一种机制是引力波辐射导致轨道不断收缩，直到伴星（在这些情形中通常为白矮星）解体$^{[8]}$。第三种机制是脉冲星进入一颗更大恒星的包层，使其解体并在脉冲星周围形成一个盘$^{[10][11]}$。由这些过程形成的盘，其质量远大于回落物质形成的盘，因此能够存活更长时间，从而允许行星形成$^{[8]}$。这些盘还富含行星形成所必需的重元素，并且盘物质的一部分会被脉冲星吸积，从而使其自转加快$^{[12]}$。此外，也可能出现另一种情形：一颗较轻的白矮星在与更大质量白矮星的相互作用中被摧毁，轻白矮星形成的碎屑盘孕育行星，而较大的白矮星最终演化为脉冲星$^{[13]}$。
\item 在某些情况下，伴星在与脉冲星的相互作用中被摧毁，但会留下一个行星尺度的残余体$^{[3]}$。这类系统被称为“黑寡妇”系统$^{[14]}$，但并不总被视为脉冲星行星的一种 $^{[15]}$。
\item 最后，还有可能是来自伴星的行星或流浪行星被脉冲星俘获 $^{[16]}$，或者脉冲星与行星原先所围绕的宿主恒星发生并合 $^{[17]}$。后一种过程会形成一个公共包层，该包层最终瓦解并形成一个盘，从中孕育行星 $^{[18]}$。
\end{itemize}
\subsubsection{影响}
不同的形成情景会对行星的组成产生直接影响：由超新星碎屑形成的行星很可能富含金属和放射性同位素$^{[16]}$，并且可能含有大量的水$^{[19]}$；由白矮星解体形成的行星则会富含碳$^{[16]}$，并可能由大量金刚石构成$^{[20]}$；而真正的白矮星碎片本身将具有极高的密度$^{[16]}$。截至 2022 年，在脉冲星周围发现的最常见行星类型是所谓的“钻石行星”，即一种极低质量的白矮星$^{[21]}$。此外，围绕脉冲星的天体还可能包括小行星、彗星以及小行星体$^{[22]}$。更具推测性的设想还包括由奇异物质构成的行星，这类行星可能比普通物质行星更靠近脉冲星运行，并有潜力发射引力波$^{[23]}$。

行星还可以与脉冲星的磁场发生相互作用，产生所谓的“阿尔芬翼”（Alfvén wings）：这是一种围绕行星形成的翼状电流结构，可向行星注入能量$^{[24]}$，并可能产生可被探测到的射电辐射$^{[25]}$。
\subsection{可观测性}
脉冲星是极其精确的“宇宙时钟”$^{[4]}$，其脉冲计时具有高度的规律性。因此，可以通过脉冲星计时变化来探测其周围非常小的天体，最小甚至可达大型小行星的尺度$^{[1]}$。在分析计时信号时，需要对地球和太阳系运动的影响、脉冲星位置估计误差，以及电磁辐射在星际介质中传播时间的变化进行校正。脉冲星会以高度规律的方式自转并逐渐减慢 $^{[4]}$；行星通过对脉冲星施加引力作用，会改变这一模式，从而在脉冲信号中产生多普勒频移$^{[26]}$。理论上，这种技术还可用于探测脉冲星行星的系外卫星$^{[27]}$。然而，该方法也存在局限：脉冲星“突变”（glitch）以及脉冲模式变化有时会模拟出行星存在的信号$^{[28]}$。

最早被发现的系外行星$^{[b]}$（1992 年由 Dale Frail 和 Aleksander Wolszczan 发现）正是围绕 PSR B1257+12 的脉冲星行星$^{[31]}$。这一发现表明，从地球上可以探测到系外行星$^{[32]}$，并由此引发了“系外行星并不罕见”的预期$^{[4]}$。截至 2016 年 $^{[33]}$，已知质量最小的系外行星（PSR B1257+12 A，仅约 $0.02,M_{\oplus}$）也是一颗脉冲星行星$^{[34]}$。

然而，由于体积极小以及脉冲星复杂而特殊的光谱特征，直接成像这类行星实际上极为困难$^{[16]}$。一种潜在的成像方式是观测行星凌日：但对脉冲星行星而言，由于脉冲星本身尺寸极小，行星发生凌日的几率非常低。同时，脉冲星的复杂光谱也使得对行星进行光谱分析变得十分困难。相比之下，行星磁场—脉冲星相互作用以及行星自身的热辐射，更有可能成为获取行星信息的途径 $^{[35]}$。

此外，脉冲星行星还被用于解释某些天文现象，例如软伽马重复暴（soft gamma repeaters）的 X 射线爆发$^{[36]}$。
\subsection{发生率}
截至 2022 年，仅有大约半打$^{[c]}$ 脉冲星行星被确认$^{[11]}$，这意味着其发生率不超过每 200 颗脉冲星中出现 1 个行星系统$^{[d][40]}$。由于对“行星”定义的差异，不同巡天项目给出的统计数量也不尽相同。大规模脉冲星行星搜索的结果同样支持这种极端稀有性$^{[41]}$。大多数行星形成情景要求其前身为质量差异很大的双星系统，且该系统必须在形成脉冲星的超新星爆发中幸存下来；而这两种条件本身都极少同时满足，因此脉冲星行星的形成是一个罕见过程$^{[3]}$。此外，行星及其轨道还必须经受住脉冲星释放的强烈高能辐射（包括 X 射线、伽马射线和高能粒子，即“脉冲星风”）$^{[6]}$。对于通过吸积而“回旋加速”的毫秒脉冲星而言，这一点尤为关键：在其作为 X 射线双星的阶段，所释放的辐射可能会蒸发行星$^{[42]}$。再者，脉冲星的可观测寿命通常只有几百万年，短于行星形成所需的时间，从而进一步限制了我们观测到它们的机会 $^{[43]}$。

基于已知的发生率估计，银河系中可能存在多达一千万颗脉冲星行星$^{[e][46]}$。目前所有已知的脉冲星行星都位于毫秒脉冲星周围$^{[1]}$，这些都是通过从伴星吸积物质而被“加速自转”的年老脉冲星。截至 2015 年，尚未发现围绕年轻脉冲星的行星$^{[47]}$；这是因为年轻脉冲星的脉冲规律性较差，计时误差更大，从而使行星探测更加困难$^{[35]}$。
\subsubsection{已确认的脉冲星行星}
\begin{figure}[ht]
\centering
\includegraphics[width=14.25cm]{./figures/3cd1c1495c4b2471.png}
\caption{} \label{fig_Pupl_2}
\end{figure}
\textbf{M62H}

M62H 是一颗位于蛇夫座的毫秒脉冲星。它位于球状星团 Messier 62 中$^{[62]}$，距离地球约 5 600 秒差距（18 000 光年$^{[63]}$。该脉冲星于 2024 年通过 MeerKAT 射电望远镜被发现$^{[62]}$。M62H 的自转周期为 3.70 毫秒，这意味着它每秒完成约 270 次自转（270 Hz$^{[64]}$。其行星伴星的最小质量为 2.5 个木星质量（$M_J$），在假设脉冲星质量为 1.4 个太阳质量（$M_\odot$）的情况下，其中值质量为 2.83$M_J$。其最小密度为 11 g/cm$^3$。若采用中值质量，则意味着其最大半径约为 48 850 千米（30 350 英里）$^{[65]}$。该行星仅需 0.133 天（3.2 小时）即可完成一次公转，其轨道半径相当于天文单位的 0.49\%，绕 M62H 运行$^{[66]}$。

\textbf{PSR B1257+12}

脉冲星 PSR B1257+12 位于室女座，距离地球约$710^{+43}_{-38}$秒差距$^{[67]}$，其拥有行星的事实于 1992 年通过阿雷西博天文台的观测得到确认$^{[68]}$。该系统包含一颗质量仅为$0.02\pm0.002$ 个地球质量的小型行星，以及两颗质量分别为$4.3\pm0.2$和$3.9\pm0.2$ 个地球质量的超级地球，假设脉冲星质量为 1.4 个太阳质量$^{[69]}$。这些行星极有可能形成于一个原行星盘$^{[1]}$，该盘可能源自伴星的部分破坏$^{[8]}$。计算机模拟表明，该系统至少在十亿年的时间尺度上是稳定的$^{[69]}$，并且系外卫星也有可能在该系统中长期存活$^{[70]}$。该系统在结构上类似于太阳系内区$^{[4]}$；这些行星绕脉冲星运行的距离与水星绕太阳的距离相当，其表面温度也可能处于类似水平$^{[71]}$。关于该系统中存在额外天体的部分报道，可能源于太阳活动造成的扰动$^{[72]}$。

\textbf{PSR J1719−1438}

一颗冥王式行星（cthonian planet）$^{[73]}$ 围绕脉冲星 PSR J1719−1438运行，其质量与木星相当，但半径不足木星的 40\% $^{[i][1]}$。该行星很可能是原伴星在脉冲星强辐射作用下被蒸发后留下的富碳残骸$^{[3]}$，因此常被称为一颗“钻石行星”$^{[j][6]}$。

\textbf{PSR B1620−26}

一颗质量约为$2.5\pm1$ 个木星质量的环双星行星$^{[75]}$，围绕位于球状星团 M4 中的双星系统 PSR B1620−26 运行$^{[4]}$。该双星系统由一颗脉冲星和一颗白矮星组成$^{[1]}$。这颗行星可能是通过引力俘获进入脉冲星轨道的，在球状星团这种高密度环境中，该过程尤为可能发生$^{[16]}$。其年龄可能高达约 126 亿年，使其成为已知最古老的行星$^{[k][76]}$。这一发现也表明，即使在低金属丰度的环境（如球状星团）中，行星同样可以形成$^{[77]}$。

\textbf{PSR J2322−2650}

PSR J2322−2650 似乎拥有一颗质量大致与木星相当的伴星。来自脉冲星的辐射可能将该伴星加热至约 2300 K；在脉冲星附近观测到的一个光源，可能正是这颗行星$^{[78]}$。该脉冲星的光度显著低于许多其他脉冲星，这或许解释了该行星为何能够存活至今$^{[79]}$。利用 JWST NIRSpec 的观测发现，这颗行星的大气层富含分子态碳（如 C$_3$、C$_2$），并存在强烈的向西风$^{[80]}$。
\subsubsection{碎屑盘与前身天体}
对脉冲星 PSR B1937+21 和 PSR J0738−4042 的计时变化，可能反映了这些脉冲星周围存在一条小行星带$^{[l]}$。此外，有学者提出，小行星或彗星与脉冲星发生碰撞，可能是快速射电暴$^{[m]}$、伽马射线暴 GRB 101225A$^{[6]}$ 以及其他类型脉冲星变异现象的成因之一$^{[84]}$。目前尚未在脉冲星周围发现已知的碎屑盘，不过有研究认为磁星 4U 0142+61 和 1E 2259+586$^{[n]}$ 可能拥有此类结构 $^{[2]}$。

白矮星—脉冲星双星系统 PSR J0348+0432 可能是一个未来有潜力形成脉冲星行星的系统$^{[86]}$。此外，也有人提出，在脉冲星 Geminga 周围存在的尘埃云，可能是行星形成的前身结构$^{[87]}$。
\subsubsection{候选体}
早期曾有一些关于脉冲星行星的报道，后来被撤回或被认为证据不足$^{[88]}$。例如，1991 年所谓在 PSR B1829−10 周围发现行星的“发现”，后来被证明只是由地球运动引起的观测伪影$^{[4]}$。自 1979 年以来，关于脉冲星 PSR B0329+54 是否存在行星的问题一直存在争议，截至 2017 年仍未得到解决$^{[89]}$。而 PSR B1828−11 已被明确证实，其表现出的现象源于磁层活动，会模拟出类似行星的信号，但实际上并不存在行星$^{[90]}$。同样，曾提出的围绕脉冲星 Geminga 的行星候选体，后来也被归因于计时噪声$^{[87]}$。
\begin{figure}[ht]
\centering
\includegraphics[width=14.25cm]{./figures/9dea5ff9cc1ac3a9.png}
\caption{} \label{fig_Pupl_3}
\end{figure}
\subsection{宜居性}
脉冲星发射的辐射谱与普通恒星截然不同：几乎不产生可见光或红外辐射，但会释放大量电离辐射$^{[46]}$以及电子–正电子对，这些粒子由脉冲星在自转过程中其强磁场所产生。此外，脉冲星诞生前遗留下来的热量、其自身辐射对两极的加热以及物质吸积过程，都会驱动热辐射和中微子的发射$^{[100]}$。电子–正电子对和 X 射线会被行星大气吸收并加热，从而引发强烈的大气逃逸，甚至可能将大气彻底剥离$^{[101]}$。若行星具备自身的磁场，可能在一定程度上缓解电子–正电子对的影响$^{[102]}$。

传统上，宜居性通过行星的平衡温度来定义，而平衡温度取决于其接收的辐射通量；若行星表面能够存在液态水，则被称为“宜居”$^{[103]}$。不过，即便外部能量输入很少，一些行星也可能在地下孕育生命$^{[104]}$。由于脉冲星体积很小，其辐射总量并不大，因此其宜居带往往会非常靠近恒星本身，以至于潮汐效应可能直接摧毁行星$^{[105]}$。此外，对于某一具体脉冲星，其实际辐射强度以及有多少辐射能够到达假想行星表面，往往难以准确判定；在已知的脉冲星行星中，只有 PSR B1257+12 系统中的行星接近宜居带$^{[106]}$，而截至 2015 年，没有任何已知的脉冲星行星被认为可能是宜居的$^{[4][39]}$。额外的热源还可能包括在导致脉冲星形成的超新星过程中生成的放射性同位素（如 钾-40）$^{[19]}$，以及近轨道行星所经历的潮汐加热$^{[107]}$。来自伴星等外部天体的辐射，也会增加行星的能量收支$^{[73]}$。
\subsection{参见}
\begin{itemize}
\item 行星列表
\item 具有原行星盘的恒星列表
\item 安德鲁·莱恩（Andrew Lyne）
\end{itemize}
\subsection{注释}\\
a.超新星爆发前就已存在并在爆发中幸存下来的行星，被称为“蝾螈（Salamander）情景”；在神话中，蝾螈被认为能够在火焰中存活。由恒星碎屑形成的行星被称为“墨农之裔（Memnonides）情景”；据罗马诗人奥维德记载，墨农之裔是由战士墨农的灰烬化成的鸟类。$^{[5]}$\\
b.早期对 HD 114762 b 和 Gamma Cephei Ab 的探测在当时被认为并不确定，因此它们不被视为最早发现的系外行星；$^{[29]}$ 此外，HD 114762 b 后来被发现实际上是一颗恒星（红矮星而非行星）。$^{[30]}$\\
c.截至 2023 年 3 月 25 日，NASA 系外行星档案库中以“PSR”为名称的天体列出了 7 颗行星 $^{[37]}$，而系外行星百科全书在相同标准下列出了 24 颗行星。$^{[38]}$\\
d.作为对比，人们认为大约四分之一到五分之一的已知白矮星（另一类恒星遗骸）拥有行星。$^{[39]}$\\
e.再作比较，银河系大约拥有 1000–4000 亿 颗恒星，$^{[44]}$ 其中大多数被认为都拥有行星。$^{[45]}$\\
f.行星按发现顺序命名，从恒星名后加一个小写字母 b 开始。在多恒星系统中，恒星在系统名后用大写字母标记，主星从 A 开始。$^{[48]}$\\
g.从地球质量换算为木星质量，应除以约 318。\\
h.半径的计算使用中位质量和最小密度，代入公式
  \[
  d=\frac{1.89813\times10^{30}, m}{(4/3)\pi r^3},~
  \]
其中$d$为密度（$g/cm{^3}$），$m$为质量（以$M_J$为单位），$r$为半径（厘米）。将厘米换算为木星半径$R_J$时，应再除以$7.1492\times10^{9}$。\\
i.有时该天体也被称为 PSR J1719-14，对应 PSR J1719-14。\\
j.其密度–质量–半径特征表明，它可能完全由金刚石构成。$^{[74]}$\\
k.另一种解释认为该行星形成于共同包层过程，这将使其年龄可能仅约 5 亿年。$^{[18]}$\\
l.在 PSR B1937+21 的情形中，最重的天体被认为质量小于地球的 1/10,000。$^{[81]}$\\
m.快速射电暴（FRB）是持续毫秒量级、起源于银河系之外的射电脉冲。$^{[82]}$ 一种理论认为，处在脉冲星磁场内轨道运行的行星会造成扰动并产生这种爆发，但目前尚无该过程的已知实例。$^{[83]}$\\
n.有资料称其名称为 1E 2259+286$^{[2]}$，但正确名称应为 1E 2259+586。$^{[85]}$\\
o.从地球质量换算为木星质量，应除以约 318。
\subsection{参考文献}
\begin{enumerate}
\item 马丁、利维奥与帕拉尼萨米 2016, 第1页。
\item 马丁、利维奥与帕拉尼萨米 2016, 第8页。
\item 马丁、利维奥与帕拉尼萨米 2016, 第4页。
\item 沃尔什赞 2015.
\item 菲尼与汉森 1993, 第371页。
\item 帕特鲁诺与卡马 2017, 第1页。
\item 马丁、利维奥与帕拉尼萨米 2016, 第2页。
\item 马丁、利维奥与帕拉尼萨米 2016, 第3页。
\item 马尔加利特与梅茨格 2017, 第2798页。
\item 平井与波德西亚德洛夫斯基 2022, 第4545页。
\item 平井与波德西亚德洛夫斯基 2022, 第4553页。
\item 欧韦尔 1992, 第668页。
\item 波德西亚德洛夫斯基、普林格尔与里斯 1991, 第783页。
\item 贝勒斯等 2011, 第1717页。
\item 莱科克与克里斯托杜鲁 2025, 第4页。
\item 内科拉·诺瓦科娃与佩特拉塞克 2017, 第1页。
\item 波德西亚德洛夫斯基、普林格尔与里斯 1991, 第784页。
\item 麦克罗伯特 2005, 第26页。
\item 帕特鲁诺与卡马 2017, 第10页。
\item 马尔加利特与梅茨格 2017, 第2800页。
\item 尼图等，2022, 第2446页。
\item 莫特兹与海瓦茨，2011, 第1页。
\item 库尔班、耿与黄 2019, 第1页。
\item 莫特兹与海瓦尔茨 2011, 第8页。
\item 莫特兹与海瓦茨 2011, 第9页。
\item 弗拉姆 1992, 第290页。
\item 刘易斯、萨克特与马德林 2008, 第156页。
\item 克尔等人 2015, 第1页。
\item 维拉斯 2016, 第1页。
\item 基弗 2019, 第1页。
\item 卡列加里、费拉兹-梅洛与米赫琴科 2006, 第381页。
\item 沃尔什赞，1994, 第542页。
\item 维拉斯，2016, 第17页。
\item 刘易斯、萨克特与马德林 2008, 第153页。
\item 内科拉·诺瓦科娃 & 佩特拉塞克 2017, 第2页。
\item 库尔班等，2024, 第2页。
\item 美国国家航空航天局2023年报告.
\item 欧洲航天局2023年报告.
\item 维拉斯与维多托 2021, 第1702页。
\item 平井与波德西亚德洛夫斯基 2022, 第4554页。
\item 莱科克与克里斯托杜鲁 2025, 第1页。
\item 米勒与汉密尔顿 2001, 第864页。
\item 米勒与汉密尔顿 2001, 第869页。
\item Stellato 2020, 第1页。
\item 卡桑等，2012, 第167页。
\item Patruno & Kama 2017, 第2页.
\item Spiewak等，2018, 第470页.
\item 国际天文学联合会.
\item 2024年NASA目录, PSR B1620-26 b.
\item NASA目录 2024, PSR J1719-1438 b。
\item NASA目录 2024, PSR J2322-2650 b。
\item 2024年EPE, PSR J1748-2021H b.
\item 2024年EPE, PSR J0636+5129 b.
\item 2024年EPE, PSR J1807-2459 A b.
\item EPE 2024, PSR B1802-07 b.
\item EPE 2024, PSR J1211-0633 b.
\item EPE 2024, PSR J0312-0921 b.
\item EPE 2024, PSR J1928+1245 b。
\item EPE 2024, PSR J1824-2452M b。
\item EPE 2024, PSR J2241-5236 b。
\item EPE 2024, PSR J1311-3430 b。
\item 弗莱肖尔等人 2024, 第1436页。
\item 奥利维拉等人 2022, 第1页。
\item 弗莱肖尔等人 2024, 第1440页。
\item 弗莱肖尔等人 2024, 第1454页。
\item 弗莱肖尔等人 2024, 第1444页。
\item 颜等，2013, 第166页。
\item 考温，1994, 第151页。
\item 沃尔什赞，2008, 第2页。
\item 多尼森，2010, 第1919页。
\item 沃尔什赞与弗雷尔，1992, 第146页。
\item 汉森、施和柯里，2009, 第387页。
\item 伊奥里奥 2021, 第1页。
\item 史密斯等 2014, 第3页。
\item 沃尔什赞 2008, 第3页。
\item 帕斯夸与阿萨夫 2014, 第1页。
\item 塞蒂亚万等 2010, 第1642页。
\item 斯皮瓦克等 2018, 第474页。
\item 斯皮瓦克等，2018年, 第476页。
\item 张，迈克尔；贝莱兹奈，玛雅；布兰特，蒂莫西·D.；罗马尼，罗杰·W.；高，彼得；贝尔茨，海莉；贝尔斯，马修；尼克松，马修·C.；宾，雅各布·L.；科马切克，萨迪厄斯·D.；科伊，布兰登·P.；傅，广伟；卢克，拉斐尔；里尔登，丹尼尔·J.；卡利，艾玛；香农，瑞安·M.；福特尼，乔纳森·J.；皮埃特，安贾莉·A.A.；科尔曼·米勒，M.；德塞特，让-米歇尔（2025）。“一颗风速极高的脉冲星行星上存在富含碳的大气层”。arXiv:2509.04558 [astro-ph.EP].
\item 尼图等，2022, 第2455页。
\item 彼得罗夫等，2015, 第457页。
\item 彼得罗夫等，2015, 第458页。
\item 希勒等，2008, 第3页。
\item 卡普兰等，2009.
\item 安东尼阿迪斯等，2013, 第448页。
\item 格里夫斯与霍兰德，2017, 第26页。
\item 沃尔什赞，1994, 第538页。
\item 斯塔罗沃伊特与罗丁 2017, 第948页。
\item 尼图等人 2022, 第2447页。
\item NASA目录 2024, PSR B0329+54 b。
\item 尼图等 2022, 第2451页。
\item EPE 2024, PSR J1555-2908 c。
\item EPE 2024, PSR B0525+21 b.
\item EPE 2024, PSR B1937+21 b。
\item EPE 2024, SGR 1806-20 b.
\item 2024年EPE, SWIFT J1756.9-2508 b.
\item 2024年EPE, PSR B0943+10 b.
\item 2024年EPE, PSR B0943+10 c.
\item 帕特鲁诺与卡马 2017, 第4–5页。
\item 帕特鲁诺与卡马 2017, 第5–6页。
\item 帕特鲁诺与卡马 2017, 第11页。
\item 帕特鲁诺与卡马2017, 第6页。
\item 斯塔门科维奇与布罗伊尔2009, 第58页。
\item 帕特鲁诺 & 卡马 2017, 第4页。
\item 帕特鲁诺 & 卡马 2017, 第7页。
\item 伊奥里奥 2021, 第5页。
\end{enumerate}
\subsubsection{资料来源}
\begin{itemize}
\item 安东尼阿迪斯，约翰；弗雷雷，保罗·C·C.；韦克斯，诺伯特；陶里斯，托马斯·M.；林奇，瑞安·S.；范克尔克维克，马滕·H.；克拉默，迈克尔；巴萨，塞斯；迪隆，维克·S.；德里贝，托马斯；赫塞尔斯，杰森·W·T.；卡斯皮，维多利亚·M.；孔德拉季耶夫，弗拉季斯拉夫·I.；兰格，诺伯特；马什，托马斯·R.；麦克劳林，莫拉·A.；彭努奇，蒂莫西·T.；兰索姆，斯科特·M.；斯特尔斯，英格丽德·H.；范·勒文，约里；韦尔比斯特，约里斯·P.W.；惠兰，大卫·G.（2013年4月26日）。“一颗位于紧凑相对论双星中的大质量脉冲星”. 科学. 340 (6131) 1233232. arXiv:1304.6875. Bibcode:2013Sci...340..448A. doi:10.1126/science.1233232. ISSN 0036-8075. PMID 23620056. S2CID 15221098.
\item 贝尔斯，M.；贝茨，S. D.；巴莱拉奥，V.；巴特，N. D. R.；布尔盖，M.；伯克-斯波洛尔，S.；达米科，N.；约翰斯顿，S.；基思，M. J.；克拉默，M.；库尔卡尼，S. R.；列文，L.；莱恩，A. G.；米利亚，S.；波森蒂，A.；斯皮特勒，L.；斯塔珀斯，B.；范斯特拉滕，W.（2011年9月23日）。毫秒脉冲双星中恒星向行星的转变. 科学. 333 (6050): 1717–1720. arXiv:1108.5201. Bibcode:2011Sci...333.1717B. doi:10.1126/science.1208890. ISSN 0036-8075. PMID 21868629. S2CID 206535504.
\item 卡列加里，N.；费拉兹-梅洛，S.；米赫琴科，T. A.（2006年4月1日）。“处于3/2平均运动共振中的两颗行星的动力学：应用于脉冲星PSR B1257+12的行星系统”. 天体力学与动力天文学. 94 (4): 381–397. 文献标识符:2006CeMDA..94..381C. DOI:10.1007/s10569-006-9002-4. ISSN 1572-9478. S2CID 123024733.
\item 卡桑，A.；库巴斯，D.；博利厄，J.-P.；多米尼克，M.；霍恩，K.；格林希尔，J.；瓦姆斯甘斯，J.；门齐斯，J.；威廉姆斯，A.；约尔根森，U. G.；乌达尔斯基，A.；贝内特，D. P.；阿尔布罗，M. D.；巴蒂斯塔，V.；布里扬，S.；考德威尔，J. A. R.；科尔，A.；库图尔，Ch；库克，K. H.；迪特斯，S.；普雷斯特，D. 多米尼斯；多纳托维奇，J.；福克，P.；希尔，K.；凯恩斯，N.；凯恩，S.；马尔凯特，J.-B.；马丁，R.；波拉德，K. R.；萨胡，K. C.；文特，C.；沃伦，D.；沃森，B.；祖布，M.；须美，T.；希曼斯基，M. K.；库比亚克，M.；波莱斯基，R.；索辛斯基，I.；乌拉奇克，K.；皮特任斯基，G.；维尔日科夫斯基，Ł（2012年1月）。“通过微引力透镜观测发现，银河系每颗恒星至少拥有一颗或更多被束缚的行星”. 自然. 481 (7380): 167–169. arXiv:1202.0903. 文献编码:2012Natur.481..167C. DOI:10.1038/nature10684. ISSN 1476-4687. PMID 22237108. S2CID 2614136.
\item 米勒，M·科尔曼；汉密尔顿，道格拉斯·P.（2001年4月）。“PSR 1257+12行星系统对孤立毫秒脉冲星的启示”. 天体物理学杂志. 550 (2): 863. arXiv:astro-ph/0012042. 文献编码:2001ApJ...550..863M. DOI:10.1086/319813. ISSN 0004-637X. S2CID 10770838.
\item 考恩，罗恩（1994年3月5日）。“围绕脉冲星运行的行星的新证据”。《科学新闻》。第145卷，第10期。《科学新闻》。第151–152页。检索日期：2023年3月23日——经盖尔学术OneFile数据库获取。
\item 唐尼森，J. R.（2010年5月）。“系外行星可能卫星的希尔稳定性：系外卫星的稳定性”. 皇家天文学会月报. 406 (3): 1918–1934. DOI:10.1111/j.1365-2966.2010.16796.x. S2CID 117784599.
\item “目录”.系外行星百科全书. 1995. 检索日期2023年3月25日
\item “目录”.系外行星百科全书. 1995. 检索日期2024年10月9日
\item 欧维尔，E. P. J. 范登 H.（1992年4月）。“脉冲星行星”. 自然. 356 (6371): 668. 文献标识码:1992Natur.356..668V. DOI:10.1038/356668b0. ISSN 1476-4687. S2CID 186241974.
\item 弗拉姆，费伊（1992年1月17日）：“天文学家是否已发现一对脉冲星行星？：新观测结果进一步支持围绕着已耗尽燃料的恒星运行的行星这一说法——并推动了寻找其成因的探索。”科学. 255 (5042): 290. doi:10.1126/science.255.5042.290. PMID 17779576.
\item 格里夫斯，J. S.；霍兰德，W. S.（2017年10月）。“中红外与亚毫米波段的杰明加脉冲星风星云”. 皇家天文学会月报：快报. 471 (1): L26 – L30. doi:10.1093/mnrasl/slx098.
\item 汉森，布拉德·M·S.；施伊，辛怡；柯里，泰恩（2009年1月）。“脉冲星行星：类地行星形成的检验案例”. 天体物理学报. 691 (1): 382–393. arXiv:0908.0736. 文献编码:2009ApJ...691..382H. DOI:10.1088/0004-637X/691/1/382. ISSN 0004-637X. S2CID 18322234.
\item 平井亮介；波德西亚德洛夫斯基，菲利普（2022年11月2日）。“中子星与双星伴星碰撞：超高速恒星、脉冲星行星、崎岖的超亮超新星以及索恩-齐特科夫天体的形成”. 皇家天文学会月报. 517 (3): 4544–4556. arXiv:2208.00915. doi:10.1093/mnras/stac3007.
\item “系外行星命名”。国际天文学联合会。检索日期：2023年9月12日。
\item 伊奥里奥，洛伦佐（2021年7月）。“经典与广义相对论倾斜进动对恒星周围中子星行星宜居性的影响”。《天文学杂志》162（2）：51。arXiv：2106.06024Bibcode2021AJ....162...51Idoi10.3847/1538-3881/ac09f8ISSN1538-3881S2CID235417162
\item 卡普兰，大卫·L.；查克拉巴蒂，迪普托；王忠祥；瓦希特，斯蒂芬妮（2009年6月）。“磁星1E 2259+586的中红外对应体”. 天体物理学杂志. 700 (1): 149–154. arXiv:0906.1604. 文献编码:2009ApJ...700..149K. DOI:10.1088/0004-637X/700/1/149. ISSN 0004-637X. S2CID 9937378.
\item 克尔，M.；约翰斯顿，S.；霍布斯，G.；香农，R. M.（2015年8月）。“年轻脉冲星周围行星形成的限制及其对超新星吸积盘的启示”。天体物理学报快报809（1）：L11。arXiv：1507.06982文献编码2015ApJ...809L..11Kdoi10.1088/2041-8205/809/1/L11ISSN2041-8205S2CID118144284
\item 基弗，弗拉维安（2019年12月1日）。“利用盖亚数据确定行星候选者HD 114762 b的质量”。天文学与天体物理学632：L9。arXiv：1910.07835文献标识码2019A&A...632L...9Kdoi10.1051/0004-6361/201936942ISSN0004-6361S2CID204743831
\item 库尔班，阿布杜沙塔尔；耿，金俊；黄永峰（2019年7月17日）。“由合并的奇异夸克星-奇异夸克行星系统产生的引力波辐射”. AIP会议论文集. 2127 (1): 020027. 文献标识码:2019AIPC.2127b0027K. DOI:10.1063/1.5117817. ISSN 0094-243X. S2CID 199118120.
\item 库尔班，阿卜杜萨塔尔；周霞；王娜；黄永峰；王宇斌；努尔马特，努里曼古丽（2024年6月1日）。“重复X射线暴：中子星与部分从行星上解体的团块之间的相互作用”。天文学与天体物理学. 686: A87.arXiv:2403.13333. Bibcode:2024A&A...686A..87K. doi:10.1051/0004-6361/202347828.
\item 莱科克，西拉斯·G·T.；克里斯托杜鲁，迪米特里斯·M.（2025年3月20日）。“关于已确认脉冲星行星的数量：六法则”. 天体物理学杂志. 982 (1): 63. 文献标识码:2025ApJ...982...63L. DOI:10.3847/1538-4357/adb1a8.
\item 刘易斯，凯伦·M.；萨克特，佩妮·D.；马德林，罗丝玛丽·A.（2008年9月）。“通过到达时间分析探测脉冲星行星卫星的可能性”。天体物理学杂志685（2）：L153。arXiv：0805.4263文献编码2008ApJ...685L.153LDOI10.1086/592743ISSN0004-637XS2CID17818202
\item 麦克罗伯特，A. M.（2005）。“追踪那个故事：M4中的脉冲星行星”. 《天空与望远镜》. 109 (1): 26. 文献标识码:2005S&T...109Q..26M.
\item 马尔加利特，本；梅茨格，布赖恩 D.（2017年3月1日）。“白矮星—中子星双星并合形成10^29克拉钻石：脉冲星行星的起源”. 皇家天文学会月报. 465 (3): 2790–2803. arXiv:1608.08636. doi:10.1093/mnras/stw2640.
\item 马丁，丽贝卡·G.；利维奥，马里奥；帕拉尼萨米，迪维亚（2016年11月）。“脉冲星行星为何罕见？”. 天体物理学杂志. 832 (2): 122. arXiv:1609.06409. 文献标识码:2016ApJ...832..122M. doi:10.3847/0004-637X/832/2/122. ISSN 0004-637X. S2CID 118490527.
\item 莫特兹，F.；海瓦茨，J.（2011年8月1日）。“行星与小天体对脉冲星风的磁耦合”. 天文学与天体物理学. 532: A21。arXiv:1106.0657. 文献标识码:2011A&A...532A..21M. DOI:10.1051/0004-6361/201116530. ISSN 0004-6361. S2CID 26955561.
\item “系外行星——美国国家航空航天局科学”，2023年6月7日，检索日期：2024年10月9日。
\item 尼图，尤莉安娜·C；基思，迈克尔·J；斯塔珀斯，本·W；莱恩，安德鲁·G；米卡利格，米切尔·B（2022年3月29日）。“对乔德雷尔银行脉冲星测时计划中800颗脉冲星周围行星伴星的搜寻”. 皇家天文学会月报. 512 (2): 2446–2459. arXiv:2203.01136. doi:10.1093/mnras/stac593.
\item 帕斯夸，安东尼奥；阿萨夫，胡代尔·A.（2014年2月25日）。“利用到达时间分析探测绕脉冲星行星运行且轨道倾斜的系外卫星的可能性”. 天文学进展. 2014e450864。文献标识码:2014AdAst2014E...6P. DOI:10.1155/2014/450864. ISSN 1687-7969.
\item 佩特罗夫，E.；约翰斯顿，S.；基恩，E. F.；范斯特拉滕，W.；贝尔斯，M.；巴尔，E. D.；巴尔斯德尔，B. R.；伯克-斯波洛尔，S.；卡莱布，M.；钱德勒，D. J.；弗林，C.；詹姆斯恩，A.；克雷默，M.；吴，C.；波森蒂，A.；斯塔珀斯，B. W.（2015年11月21日）。快速射电暴领域综述：可重复性的限制. 皇家天文学会月报. 454 (1): 457–462. arXiv:1508.04884. doi:10.1093/mnras/stv1953.
\item “表格”.美国国家航空航天局系外行星档案. 检索日期2023年3月25日
\item 内科拉·诺瓦科娃，J.；佩特拉塞克，T.（2017年9月1日）。脉冲星行星表征的可行性和优势。2017年欧洲行星科学大会。第623页。EPSC2017–623。文献标识码:2017EPSC...11..623N.
\item 奥利维拉，R. A. P.；奥尔托拉尼，S.；巴布伊，B.；克尔伯，L. O.；马亚，F. F. S.；比卡，E.；卡西西，S.；索萨，S. O.；佩雷斯-比列加斯，A.（2022年1月）。“来自OGLE-IV成员RR Lyrae恒星在六个银心球状星团中的精确距离”。天文学与天体物理学. 657: A123.arXiv:2110.13943. Bibcode:2022A&A...657A.123O. doi:10.1051/0004-6361/202141596.
\item 帕特鲁诺，A.；卡马，M.（2017年12月1日）。“中子星行星：大气过程与辐射”. 天文学与天体物理学. 608：A147。arXiv:1705.07688. 文献标识码:2017A&A...608A.147P. DOI:10.1051/0004-6361/201731102. ISSN 0004-6361. S2CID 119191976.
\item 菲尼，E. S.；汉森，B. M. S.（1993年1月）。
\item 波德西亚德洛夫斯基，Ph；普林格尔，J. E.；里斯，M. J.（1991年8月）。“围绕PSR1829 – 10运行的行星的起源”。自然352（6338）：783–784。天文学编码：1991Natur.352..783PDOI10.1038/352783a0ISSN1476-4687S2CID4235775
\item 塞蒂亚万，约翰尼；克莱门特，雷纳尔·J.；亨宁，托马斯；里克斯，汉斯-瓦尔特；罗绍，博伊克；罗德曼，延斯；舒尔策-哈图恩，蒂姆（2010年12月17日）。“一颗围绕着贫金属外星恒星运行的巨行星”. 科学. 330 (6011): 1642–1644. arXiv:1011.6376. Bibcode:2010Sci...330.1642S. doi:10.1126/science.1193342. ISSN 0036-8075. PMID 21097905. S2CID 657925.
\item 希勒，安迪；坎尼夫，约翰；沃森，布鲁诺；涅斯特罗耶夫，维塔利；布朗，迈克尔；安德森，托本；恩马克，安妮塔；林德，彼得（2008年4月22日）。安德森，托本·E.（编）。“高时间分辨率天体物理学与极大望远镜：选择哪个波长？”. 极大望远镜：哪些波长？阿恩·阿尔德贝格退休研讨会. 6986. SPIE：94–102. 文献标识码:2008SPIE.6986E..0AS. DOI:10.1117/12.801261. S2CID 120761231.
\item 史密斯，R. F.；埃格特，J. H.；让洛兹，R.；达菲，T. S.；布劳恩，D. G.；帕特森，J. R.；鲁德，R. E.；比纳，J.；拉齐基，A. E.；哈姆扎，A. V.；王，J.；布劳恩，T.；本尼迪克特，L. X.；塞利尔斯，P. M.；柯林斯，G. W.（2014年7月）。“金刚石的斜坡压缩至五太帕”. 自然. 511 (7509): 330–333. 文献标识码:2014Natur.511..330S. DOI:10.1038/nature13526. ISSN 1476-4687. PMID 25030170. S2CID 4389771.
\item 斯皮瓦克，R；贝尔斯，M；巴尔，E D；巴特，N D R；布尔盖，M；卡梅伦，A D；钱伯斯，D J；弗林，C M L；詹姆斯恩，A；约翰斯顿，S；基思，M J；克雷默，M；库尔卡尼，S R；列文，L；莱恩，A G；莫雷洛，V；吴，C；波森蒂，A；拉维，V；斯塔珀斯，B W；范斯特拉滕，W；蒂布尔齐，C（2018年3月21日）。“PSR J2322−2650——一颗拥有行星质量伴星的低光度毫秒脉冲星”。《皇家天文学会月报》475（1）：469–477。arXiv：1712.04445doi10.1093/mnras/stx3157
斯塔罗沃伊特，E. D.；罗京，A. E.（2017年11月1日）。“关于脉冲星PSR B0329+54周围行星的存在”. 天文学报告. 61 (11): 948–953. arXiv:1710.01153. Bibcode:2017ARep...61..948S. doi:10.1134/S1063772917110063. ISSN 1562-6881. S2CID 255206063.
斯特拉托，朱迪（2020）。“银河系与扁豆”。科学视野. 43 (6): 44–49. doi:10.1080/08872376.2020.12291320. ISSN 0887-2376. JSTOR 27048035.
维拉斯，迪米特里；维多托，阿琳·A（2021年7月15日）。“主序后恒星周围行星磁层的演化”. 皇家天文学会月报. 506 (2): 1697–1703. arXiv:2106.10293. doi:10.1093/mnras/stab1772.
维拉斯，迪米特里（2016年2月）。“主序后行星系统的演化”。皇家学会开放科学3（2）150571。arXiv：1601.05419Bibcode2016RSOS....350571Vdoi10.1098/rsos.150571ISSN2054-5703PMC4785977PMID26998326
斯塔门科维奇，弗拉达；布罗伊尔，多丽丝（2009年2月1日）。“专题：2008年9月1日至3日于瑞士纳沙泰尔举行的第八届欧洲天体生物学研讨会摘要”. 生命起源与生物圈演化. 39 (1): 1–89. DOI:10.1007/s11084-008-9155-0. ISSN 1573-0875. PMID 19184520. S2CID 37433981.
弗莱肖尔，L；科龙吉乌，A；斯塔珀斯，B W；弗雷雷，P C C；里多尔菲，A；阿巴特，F；兰索姆，S M；波森蒂，A；帕德马纳布，P V；巴拉克里希南，V；克拉默，M；文卡特拉曼·克里希南，V；张，L；贝尔斯，M；巴尔，E D；布赫纳，S；陈，W（2024年4月13日）。“M62中脉冲星的发现与观测时间”. 皇家天文学会月报. 530 (2): 1436–1456. arXiv:2403.12137. Bibcode:2024MNRAS.530.1436V. doi:10.1093/mnras/stae816.
沃尔什赞，亚历山大（1994年4月22日）。“证实围绕毫秒脉冲星PSR B1257+12运行的类地行星”。科学. 264 (5158): 538–542. 天文学文献标识符:1994Sci...264..538W. DOI:10.1126/science.264.5158.538. PMID 17732735. S2CID 19621191.
沃尔什赞，A（2008年8月）。“中子星行星研究十五年”. 物理学期刊. T130 014005. 文献标识码:2008PhST..130a4005W. DOI:10.1088/0031-8949/2008/T130/014005. S2CID 122989232.
沃尔什赞，A.；弗雷尔，D. A.（1992年1月）。“毫秒脉冲星PSR1257+12周围的一个行星系统”。自然355（6356）：145–147。天文学编码：1992Natur.355..145WDOI10.1038/355145a0ISSN1476-4687S2CID4260368
沃尔什赞，亚历山大（2015）。“脉冲星行星”. 天体生物学百科全书。施普林格。第2089–2092. 文献标识码:2015enas.book.2089W. DOI:10.1007/978-3-662-44185-5_1309. ISBN 978-3-662-44184-8.
严，珍；沈志强；袁建平；王娜；罗特曼，赫尔格；阿莱夫，沃尔特（2013年7月21日）。“具有行星系统的脉冲星PSR B1257+12的甚长基线干涉测量天体测量”. 皇家天文学会月报. 433 (1): 162–169. DOI:10.1093/mnras/stt712.
\end{itemize}